\section{Reproducibility}\label{sec:repro}

The source code, the data scripts and the computational records are in the repository of the studio, and the engine release that produced every number is \engineRelease{}. Every number of Section~\ref{sec:results} and of the abstract is read from the files written by one script, \texttt{manuscripts/flicker-bounded-sequences/data/measure.py}, which writes the measurement file \texttt{flicker.json}, the macro file \texttt{flicker\_numbers.tex} and the five table files that this paper includes; no number of the results is typed in the text. The script is run from the repository root with the environment of the pipeline,
\begin{center}\small\texttt{.venv-pipeline/Scripts/python.exe manuscripts/flicker-bounded-sequences/data/measure.py}\end{center}
with the defaults of \nSeedsTwo{} seeds on the acceptance task, the sort and the ablation, and \longSeeds{} seeds on the long moves; it refuses fewer than \winMinSeeds{} seeds. The figures are drawn from \texttt{flicker.json} by \texttt{manuscripts/flicker-bounded-sequences/figures/make\_figs.py}, and the document is built by three passes of \texttt{pdflatex} from the folder \texttt{tex}. Every method is deterministic given its seed through the portable generator of the engine, so a rerun reproduces the numbers up to the floating point of the machine's FFT; the seconds differ between runs. The sequence is the module \texttt{b2\_flicker\_bounded.py} and the two published sequences are options of the module \texttt{m16\_rearrangement.py} of the engine, both under \texttt{data-pipeline/pipeline/model/holo/methods}, with the tests that guard them in the repository.
